%% ─────────────────────────────────────────────────────────────────────────────
%% Capítulo 5 — Resultados y Evaluación
%% ─────────────────────────────────────────────────────────────────────────────

\chapter{Resultados y Evaluación}
\label{cap:resultados}

% Párrafo introductorio: presentar la estructura del capítulo y el conjunto
% de datos usado para la evaluación (temporada 2023, circuitos seleccionados).
% Mencionar el entorno de ejecución (hardware, SO).

\todo{Escribir párrafo de apertura del capítulo.}

% ─────────────────────────────────────────────────────────────────────────────
\section{Configuración experimental}
\label{sec:config_experimental}

% Guía:
% - Datos: temporada F1 2023, N circuitos, M pilotos.
% - Hardware: CPU, RAM, GPU (si aplica para entrenamiento Transformer).
% - Hiperparámetros: context_laps=40, d_model=384, n_layers=8, n_heads=8,
%   MC_TOP_K=3, n_sim=100, risk_lambda=0.15.
% - División train/test: por sesión (carreras = test, FP = train).
%   Justificación de la separación temporal para series temporales.

\todo{Escribir sección de configuración experimental.}

% ─────────────────────────────────────────────────────────────────────────────
\section{Evaluación de los modelos de predicción de tiempos de vuelta}
\label{sec:eval_modelos_pace}

% Guía:
% - Métricas: MAE, RMSE en segundos (tras desnormalización).
% - Comparativa: modelo específico del piloto vs. modelo global.
% - Tabla de resultados por piloto (top 10 pilotos del campeonato).
% - Tabla de resultados por circuito (street circuits vs. permanent circuits).
% - Análisis de curvas de degradación predichas vs. reales (figuras).
% - Discusión: cuándo falla el modelo (vueltas de pit-out, SC, vueltas cortas).

\todo{Escribir sección de evaluación de modelos de pace.}

% Tabla de resultados por piloto:
% \begin{table}[H]
% \centering
% \caption{MAE y RMSE de los modelos Transformer por piloto sobre el conjunto de test
%          (temporada 2023, carreras).}
% \label{tab:resultados_transformer}
% \begin{tabular}{lrrrr}
%     \toprule
%     \textbf{Piloto} & \textbf{Driver ID} & \textbf{Vueltas entren.}
%                     & \textbf{MAE (s)} & \textbf{RMSE (s)} \\
%     \midrule
%     Max Verstappen  & 1   & 312 & -- & -- \\
%     Sergio Pérez    & 11  & 298 & -- & -- \\
%     Lewis Hamilton  & 44  & 287 & -- & -- \\
%     \vdots          &     &     &    &    \\
%     \textit{Global} & --  & --  & -- & -- \\
%     \bottomrule
% \end{tabular}
% \end{table}

% ─────────────────────────────────────────────────────────────────────────────
\section{Validación del motor de estrategia}
\label{sec:validacion_motor}

% Guía:
% - Para un conjunto de carreras históricas de 2023:
%   (1) Generar las top-5 estrategias con RaceScope.
%   (2) Identificar la estrategia real usada por el ganador.
%   (3) Medir en qué posición del ranking aparece la estrategia real.
% - Tabla: circuito, estrategia real, posición en el ranking, expected_time.
% - Discusión: ¿en qué carreras se equivoca RaceScope y por qué?

\todo{Escribir sección de validación del motor de estrategia.}

% ─────────────────────────────────────────────────────────────────────────────
\section{Análisis de sensibilidad}
\label{sec:analisis_sensibilidad}

% Guía:
% - Variación de \lambda (0.0, 0.05, 0.10, 0.15, 0.20, 0.25).
% - Variación de pit_loss (20s, 22.5s, 25s, 28s).
% - Variación de P(SC) (0.10, 0.20, 0.30).
% - Variación de MC_TOP_K (3, 5, 8, 10) vs. n_sim (50, 100, 200, 500).

\todo{Escribir sección de análisis de sensibilidad.}

% ─────────────────────────────────────────────────────────────────────────────
\section{Rendimiento del sistema}
\label{sec:rendimiento_sistema}

% Guía:
% - Resultados del benchmark_strategy.py: latencias cold, warm, hot.
% - Tamaño de los modelos (.joblib): tamaño medio por piloto, total.
% - Tamaño de la caché de curvas de pace.
% - ¿Cumple el sistema el requisito de latencia < 2s en régimen caliente?

\todo{Escribir sección de rendimiento del sistema.}

% ─────────────────────────────────────────────────────────────────────────────
\section{Discusión}
\label{sec:discusion}

% Guía:
% - ¿Qué aporta RaceScope frente a trabajos relacionados
%  ?
% - Limitaciones actuales del sistema.
% - ¿En qué condiciones es más/menos fiable?

\todo{Escribir sección de discusión.}
